\section{Methodology}\label{sec:method}

BaiZe is a decoder-only (Llama-style) model built and trained from scratch with the Megatron-SWIFT framework (v4.5.3) on top of Megatron-Core.
This section details (i) the architecture and its selection rationale, (ii) the training framework and $\mu$P transfer, (iii) the WSD schedule, and (iv) the budget-constrained search protocol.

\subsection{Model Architecture}\label{sec:arch}
BaiZe follows the MiniCPM5-1B configuration: a deep-and-thin decoder-only transformer with Pre-RMSNorm and SwiGLU activations.
Grouped-query attention (GQA) shares $2$ KV heads across $16$ query heads, and word embeddings are \emph{untied} from the output head.
The full configuration is listed in Table~\ref{tab:config}; the model totals $\approx$1.04B parameters.

\begin{table}[htbp]
\centering
\caption{BaiZe model configuration.}
\label{tab:config}
\begin{tabular}{lc}
\toprule
\textbf{Hyper-parameter} & \textbf{Value} \\
\midrule
Hidden size & 1536 \\
Intermediate size & 4608 \\
Transformer layers & 24 \\
Attention heads (Q) & 16 \\
KV heads (GQA) & 2 \\
Head dimension & 96 \\
Vocabulary size & 130\,560 \\
Max position embeddings & 131\,072 \\
RoPE base ($\theta$) & 5\,000\,000 \\
RMSNorm $\varepsilon$ & $1\times10^{-6}$ \\
Tie word embeddings & false \\
Activation & SwiGLU \\
Precision & BF16 \\
Parameters & $\approx$1.04B \\
\bottomrule
\end{tabular}
\end{table}

\subsection{Architecture Selection Rationale}\label{sec:arch-rationale}
The architecture is chosen on three grounds.
\textbf{(i) A proven 1B skeleton.} At the 1B scale, the deep-and-thin GQA design is a mature, throughput-optimised configuration with stable training dynamics (Sections~\ref{sec:search}--\ref{sec:arch2b}); Megatron-Core ships no released 1B hybrid (attention+SSM) or Nemotron-style variant, making it the only instantly executable 1B architecture under our framework.
\textbf{(ii) Controlled 2B evidence.} Rather than rely on availability alone, we ran a matched from-scratch comparison at 2B between MiniCPM5-2B (dense) and a Nemotron-H-style hybrid, Mamba2-hybrid (56 layers interleaving 4 attention layers with SSM/MLP)~\cite{dao2024transformers,nvidia2025nemotronh}. Both models were trained for 1000 steps at sequence length 4096 with identical data, seed, and a 6-GPU (TP1/DP6) footprint. The hybrid reaches a \emph{lower} final loss (4.6846 vs.\ 4.8065) with \emph{fewer} parameters (2.220B vs.\ 2.512B, $-11.6\%$) and faster inference (prefill $+10.5\%$, decode $\times10.3$), at a $\sim$24\% training-throughput cost (Table~\ref{tab:arch2b}).
\textbf{(iii) A selection map.} Together these results define the architecture roadmap: the dense GQA design remains the right backbone for the budget-constrained 1B agent built here, because it maximises training throughput and has no 1B hybrid counterpart; Mamba2-hybrid is the preferred candidate when scaling the family to 2B, where inference---especially decode---efficiency dominates total cost.

\subsection{Training Framework and $\mu$P Transfer}\label{sec:framework}
We train with the Megatron-SWIFT pre-training entry point (\texttt{megatron pt}) using the \texttt{llama}/\texttt{minicpm5} model bridge.
Training is run under maximal-update parameterization ($\mu$P)~\cite{yang2022tensor}, which scales weight initialisation and learning rates according to hidden width so that hyper-parameters tuned in short, cheap runs transfer to the full configuration.
We exploit this by running all architecture and hyper-parameter sweeps at a reduced iteration count (80--160 steps) and treating the resulting ranking as predictive of the full-model optimum.
Fused attention (\texttt{attention\_backend fused}), selective activation recomputation, and cross-entropy loss fusion are enabled to raise throughput; no tensor or pipeline parallelism is required ($\texttt{tensor\_model\_parallel\_size}=1$), keeping the configuration simple and portable.

\subsection{Warmup--Stable--Decay Schedule}\label{sec:wsd}
BaiZe uses a WSD learning-rate schedule with three phases: a short linear \emph{warmup} ($10$ iterations), a long constant-rate \emph{stable} phase, and an exponential \emph{decay} to a minimum rate of $0.1\times$ the stable rate.
The decay window $\delta$ (the \emph{WSD ratio}) is itself a swept hyper-parameter.
Throughout, the global batch size is $1024$ with sequence length $4096$ and packing enabled, yielding $\approx$4\,M tokens per step.
Initially the stable learning rate is $6\times10^{-4}$ with $\delta=10\%$; the search in Section~\ref{sec:search} perturbs these two quantities together with architectural variables.

\subsection{Budget-Constrained Search Protocol}\label{sec:search}
Because a full grid is infeasible under a fixed GPU budget, we adopt a staged, single-variable-at-a-time protocol:

\begin{itemize}
\item \textbf{S1 (baseline).} Reproduce the reference MiniCPM5-1B configuration to establish a loss/throughput baseline.
\item \textbf{S2 (architecture sweep).} Perturb one architectural variable at a time---layer count ($24\rightarrow28$), KV-head count ($2\rightarrow4$), and embedding tying (untied$\rightarrow$tied)---at short horizon (80 steps).
\item \textbf{S3 (schedule sweep).} Perturb the stable learning rate ($6\times10^{-4}$ vs.\ $3\times10^{-4}$ vs.\ $1\times10^{-3}$) and the WSD ratio ($10\%$ vs.\ $20\%$).
\item \textbf{S4 (verification).} Retrain the winning configuration from scratch over a longer horizon (160 steps) to confirm convergence and reproducibility.
\end{itemize}

At every stage only a single variable changes relative to the running best, isolating each factor's contribution while keeping total GPU-hours within budget.
We also validated the architectural search space up front: hybrid attention+SSM (Hymba-style) and Nemotron alternatives were found to have no released 1B-scale variants in Megatron-Core, so the Llama-style GQA baseline remains the only executable 1B architecture.